\chapter{Confronto consolidato dei tre task}
\label{chap:results}

\paragraph{Riassunto del capitolo.} Il capitolo aggrega i risultati dei
tre task in un'unica tabella di confronto, ordina i task per difficoltà
empirica (Task 2 $\prec$ Task 1 $\prec$ Task 3) e distilla tre lezioni
trasversali: \emph{(i)} la giusta strategia di cross-validation vale
più del giusto modello, \emph{(ii)} il GridSearchCV non porta
miglioramenti misurabili senza un intervallo di confidenza, \emph{(iii)}
il feature engineering minimale e fisicamente motivato batte ensemble
ciechi più sofisticati. Verifichiamo infine la coerenza fra i numeri
dei notebook esplorativi e quelli prodotti dal codice di produzione.

\section{Tabella consolidata}

\begin{table}[H]
  \centering
  \small
  \caption{Modelli vincenti e metriche-chiave per ciascun task (run di
  default, seed 42).}
  \label{tab:results-overall}
  \begin{tabularx}{\linewidth}{@{}p{2.6cm} p{3.2cm} X@{}}
    \toprule
    \textbf{Task} & \textbf{Miglior modello} & \textbf{Metrica chiave} \\
    \midrule
    Task 1 -- Leave-One-Out & KNN ($k{=}10$, weight \texttt{distance}) & $R^{2} = 0.9909$ \\
    Task 2 -- Young / Old & SVM-RBF ($C{=}10$, $\gamma{=}$\texttt{scale}) & Acc. $= 1.0000$ sulle 8 curve del hold-out $(Aging{=}4, SOC{=}3)$; degrada al confine $\sim 0.50$ per coppie come $(3, 4)$ \\
    Task 3 -- Aging Interpolation & KNN ($k{=}10$, weight \texttt{distance}) & $R^{2} = 0.9862$ \\
    \bottomrule
  \end{tabularx}
\end{table}

\section{Difficoltà relativa}

Empiricamente, il \emph{ranking} di difficoltà è:

\[
\text{Task 2} \;\prec\; \text{Task 1} \;\prec\; \text{Task 3}.
\]

\begin{itemize}
  \item \textbf{Task 2} è il più semplice agli estremi del dominio
        ($Aging \in \{0, 4\}$): le 8 curve di Nyquist della coppia
        sono nettamente distinguibili fra young e old e quasi qualunque
        modello satura. Diventa invece un problema vero \emph{al
        confine}, attorno a $Aging \in \{2, 3\}$ --
        \cref{tab:task2-cross-pair} mostra accuracy che oscillano
        fra $0.50$ e $1.00$ a seconda del SOC scelto.
  \item \textbf{Task 1} è di difficoltà intermedia: la (Aging,
        Temperatura) esclusa è circondata, su almeno una delle due
        dimensioni, da combinazioni viste durante l'addestramento. Il
        Gradient Boosting cattura i residui non-lineari su entrambi gli
        assi.
  \item \textbf{Task 3} è il più difficile perché un \emph{intero stadio
        di aging} non è osservato. I tree-based modelli \emph{non
        estrapolano} (le loro foglie non possono predire al di fuori
        dell'inviluppo dei target visti); soltanto interpolatori
        \emph{distance-based} come KNN tengono il ritmo. La performance
        degrada inoltre ai bordi della finestra termica.
\end{itemize}

\section{Lezioni trasversali}

\subsection{La giusta CV vale più del modello}

In tutti e tre i task, la differenza fra usare una CV \emph{leakage-free}
(\acrshort{logo} o \texttt{StratifiedGroupKFold}) e una CV row-wise vale
più di una settimana di tuning di iperparametri. I numeri di
\cref{tab:task1-benchmark,tab:task2-benchmark,tab:task3-benchmark} sono
\emph{realistici} grazie alla scelta di splitter coerenti con il task.

\subsection{Tuning $\neq$ miglioramento}

In tutti e tre i task il GridSearchCV \emph{non migliora} la baseline. È
un risultato controintuitivo ma istruttivo: con dataset piccoli
(\num{9805} righe), la varianza fra fold di CV domina la varianza fra
configurazioni di iperparametri. Il valore del tuning sta nella
\emph{difesa} della scelta -- mostrare che i default non sono fortunati
-- più che nel guadagno marginale.

\subsection{Feature engineering minimale, fisicamente motivato}

Le feature più predittive in tutti i task sono quelle che riflettono
direttamente la fisica:
\begin{itemize}
  \item \texttt{inv\_Temp} (Arrhenius) batte \texttt{Temperature} pura;
  \item \texttt{log\_Freq} cattura la separazione dei processi
        elettrochimici su decadi;
  \item le interazioni \texttt{X\_x\_Temp} sono la rappresentazione più
        compatta della dipendenza incrociata fra invecchiamento e
        temperatura osservata nei dati.
\end{itemize}

Aggiungere feature più complesse (es.~modelli equivalenti di circuito
fitati) avrebbe portato un guadagno marginale a fronte di un costo
significativo in interpretabilità.

\section{Coerenza fra notebook e produzione}

Una verifica importante è che i numeri prodotti dai notebook esplorativi
in \codefile{notebooks/} coincidano con quelli prodotti dal codice di
produzione in \codefile{src/}. Per i tre task lo abbiamo verificato:
\codefile{run\_leave\_one\_out}, \codefile{run\_classification} e
\codefile{run\_aging\_interpolation} restituiscono le stesse metriche con
la stessa configurazione, perché \emph{usano gli stessi moduli} di
feature engineering, scaling e modelli. La differenza fra notebook e
produzione è solo l'ergonomia: i notebook hanno markdown, plot Plotly e
celle interattive; la produzione restituisce dataclass tipizzati che
attraversano l'API.
